\chapter{Dataset}
\label{ch:dataset}

This analysis uses the Taiwanese Bankruptcy Dataset, a well-established benchmark
compiled from the Taiwan Economic Journal \footnote{\url{https://www.tejwin.com}} from 1999 to 2009 and made accessible
through the \textcite{uci_taiwan}. Corporate bankruptcy was defined based on the
rules of the Taiwan Stock Exchange.\footnote{\url{https://twse-regulation.twse.com.tw/EN/law/DOC01_print.aspx?FLCODE=FL007304&FLNO=50-1}} 
It contains $6{,}819$ companies, each described by $95$ financial attributes covering standard business performance metrics such as 
Return on Assets (ROA), the Debt Ratio, Cash Flow Rate or Earnings per Share (EPS). Outcome is simply measured by one binary variable 
of either 0 = healthy or 1 = bankrupt.

\begin{table}[htbp]
  \centering
  \caption{Composition of the Taiwanese Bankruptcy Dataset.}
  \label{tab:dataset}
  \begin{tabular}{lrr}
    \toprule
    \textbf{Class} & \textbf{Firms} & \textbf{Share} \\
    \midrule
    Healthy (0)   & $6{,}599$ & $96.77\%$ \\
    Bankrupt (1)  & $220$     & $3.23\%$  \\
    \midrule
    Total         & $6{,}819$ & $100.00\%$ \\
    \bottomrule
  \end{tabular}
  \floatfoot{Source: \textcite{uci_taiwan}.}
\end{table}

Given this setup, two problems arise that shape the entire analysis. First one being the obvious severe class imbalance 
already visible in Table~\ref{tab:dataset}. By the nature of the problem, bankruptcies are rare events, here making up only 3.23\% 
of all observations. A naive classifier trained on the full data could reach very high accuracies just by labeling almost every company as 
healthy, but it would be predictively worthless. It gets the 96.77\% of healthy firms right while missing the very cases we care about, 
the 3.23\% that actually go bankrupt. 

The second problem is the redundancy that comes with this high dimensionality. Many of the 95 features are just algebraic transformations 
or close relatives of one another. \emph{Return on Assets} for example has ROA(A) before interest and after tax and ROA(B) before interest and 
depreciation after tax, which differ only in the treatment of depreciation. Another is \emph{Operating Profit Per Share} and \emph{Operating profit/Paid-in capital}.
Both put the same operating profit on top and divide it by closely related bases, shares outstanding in one case and the capital paid in for those shares
in the other, so they move closely together too. As a result, distinct columns can carry very similar information. 
Both of these problems need to be understood and addressed in order to build an accurate and interpretable
model.

% redundancy vielleicht doch multicolinearity
